% =============================================================================
% KAPITEL 8: FAZIT UND AUSBLICK
% =============================================================================

\chapter{Fazit und Ausblick}\label{chap:fazit}

% =============================================================================
\section{Fazit}
\label{sec:fazit}
% =============================================================================

Diese Arbeit hat untersucht, ob und in welchem Umfang die Integration von
Expected-Goals-Daten (xG) die Vorhersagegenauigkeit von Machine-Learning-Modellen
für Bundesliga-Spielergebnisse verbessert. Ausgangspunkt war die Beobachtung,
dass traditionelle Spielstatistiken -- Tore, Schüsse, Punkte -- zufallsbedingte
Schwankungen nicht herausrechnen, während xG die Qualität von Torchancen
unabhängig vom tatsächlichen Torerfolg misst und damit einen stabileren
Leistungsindikator darstellt \cite{AnzerBauer2021,Wheatcroft2022}.

\textbf{Zur zentralen Forschungsfrage:} \emph{Verbessert die Integration
von Expected-Goals-Daten die Vorhersagegenauigkeit signifikant?}

Die Antwort ist \textbf{ja} -- mit der wichtigen methodischen Einschränkung,
dass der Effekt moderat, aber konsistent ist. In allen vier paarweisen
Vergleichen zwischen xG- und Nicht-xG-Varianten zeigt die xG-Integration
eine Verbesserung des F1-macro-Scores, mit einem mittleren Gewinn von
$+0{,}041$ Punkten (von durchschnittlich 0,432 auf 0,473). Dieser Befund
ist robust: Er tritt unabhängig davon auf, ob der Datensatz bereinigt,
feature-selektiert oder unverändert verwendet wird. Insbesondere die
Expected-Goals-Daten des Auswärtsteams erweisen sich laut
Feature-Importance-Analyse als bedeutende Prädiktoren, was darauf hindeutet,
dass xG-Metriken Information kodieren, die durch die EWA-basierten
Torfeatures allein nicht vollständig erfasst wird.

Das beste Einzelmodell -- ein Random Forest auf Variante~6 (\emph{Mit xG,
Raw, Starke Features}) -- erreicht einen F1-macro von 0,479 und eine
Accuracy von 0,499. Diese Werte liegen im unteren Bereich des in der
Literatur für Pre-Match-Vorhersagen beschriebenen realistischen Korridors
von 0,45--0,55 \cite{BunkerThabtah2019,HubacekSourekZelezny2019} und sind
damit als methodisch solide, aber nicht ungewöhnlich einzuordnen. Der
erhebliche Abstand zu Post-Match-Modellen -- die mit matchweisem xG
Accuracies von über 65\,\% erzielen \cite{ForcherEtAl2025} -- erklärt sich
aus dem grundlegend anderen Informationsstand: Pre-Match-Vorhersagen müssen
ohne Kenntnis des Spielverlaufs auskommen.

Neben der Beantwortung der Forschungsfrage liefert die Arbeit drei
methodische Beiträge: Erstens eine vollständige, reproduzierbare
ML-Pipeline zur Bundesliga-Spielvorhersage mit konsequenter Data-Leakage-Prävention
durch den Lag-1-Ansatz und die Shift-Operation bei der EWA-Berechnung.
Zweitens den Nachweis, dass gezielt konstruierte Unentschieden-Features
(\texttt{combined\_draw\_tendency}, \texttt{abs\_xG\_diff}) in Verbindung
mit gewichtetem Training zu überdurchschnittlichen F1-Draw-Werten von
0,305 bis 0,385 führen -- einem in der Literatur typischerweise
schwierig zu erreichenden Ergebnis \cite{HeGarcia2009}. Drittens eine
systematische Evaluation von zwölf Datensatzvarianten, die zeigt, dass
PCA bei baumbasierten Modellen keinen Mehrwert bietet und IQR-Bereinigung
bei EWA-geglätteten Features überflüssig ist.

% =============================================================================
\section{Ausblick}
\label{sec:ausblick}
% =============================================================================

Die Ergebnisse dieser Arbeit öffnen mehrere Richtungen für zukünftige
Forschung:

\textbf{Matchweises xG:} Der deutlichste Hebel zur Verbesserung der
Vorhersagegüte liegt in der Verwendung spielphasenweiser xG-Daten.
Wenn für jedes Spiel der aktuelle xG-Durchschnitt der letzten $n$ Partien
bekannt ist -- analog zur EWA-Berechnung für Tore -- würde die xG-Information
denselben Formcharakter erhalten wie die übrigen Features. Anbieter wie
Statsbomb oder Wyscout stellen solche Daten kommerziell zur Verfügung;
im akademischen Bereich könnten öffentliche Tracking-Datensätze genutzt
werden \cite{PappalardoEtAl2019}.

\textbf{Temporal-konsistentes Evaluationsdesign:} Eine Aufteilung nach
Saisons -- Training auf 2016/17--2022/23, Validierung auf 2023/24,
Test auf 2024/25--2025/26 -- würde die zeitliche Kausalität sicherstellen
und die Übertragbarkeit der Modelle auf reale Prognoseszenarien besser
abbilden.

\textbf{Erweiterung des Feature-Sets:} Weitere Einflussfaktoren wie
Marktwerte der Kader \cite{RibeiroEtAl2025}, Reisedistanz vor Auswärtsspielen
oder die Anzahl der Tage seit dem letzten Spiel (Müdigkeitseffekt) könnten
das Modell um Kontextinformationen ergänzen, die im aktuellen Feature-Set
fehlen.

\textbf{Ensemble-Kombinationen:} Die geringen Leistungsunterschiede zwischen
Random Forest, LightGBM und XGBoost legen nahe, dass ein Voting-Ensemble
aus allen drei Modellen -- wie von \textcite{MillsEtAl2024} für ähnliche
Aufgaben gezeigt -- die Generalisierungsleistung weiter stabilisieren könnte.

\textbf{Erweiterte Ligazusammensetzung:} Die Übertragbarkeit der Befunde
auf andere Ligen -- insbesondere solche mit abweichenden Torquoten oder
anderen Heimvorteils-Charakteristika -- wäre eine naheliegende Erweiterung,
die die Generalisierbarkeit des Ansatzes validiert.

Zusammenfassend zeigt diese Arbeit, dass Expected Goals auch in einem
saisonalen Pre-Match-Format einen messbaren und konsistenten prädiktiven
Beitrag leisten. Gleichzeitig verdeutlicht sie, wo die fundamentalen
Grenzen der Vorhersagbarkeit von Fußballergebnissen liegen: nicht
primär im Modell, sondern in der Natur des Spiels selbst -- einem
Low-Scoring-Sport, in dem Zufall und nicht beobachtbare Faktoren
zwangsläufig einen erheblichen Anteil der Ergebnisvarianz erklären
\cite{DobsonGoddard2011,Sumpter2016}.